% © 2026 Ing. David Jaroš — CC BY-NC-ND 4.0
%
% This work is licensed under a Creative Commons Attribution-NonCommercial-NoDerivatives
% 4.0 International License (CC BY-NC-ND 4.0).
%
% B_base = N_eff^{3/2} — Approach A2: Rigorous motivation of the exponent
% Track: CORE — α derivation — v58 baseline
% Date: 2026-03-07

\ifdefined\INCLUDEMODE
  % Being included in another document — skip preamble
\else
  \documentclass[11pt,a4paper]{article}
  \usepackage{amsmath,amssymb,amsthm}
  \usepackage{hyperref}

  \newtheorem{theorem}{Theorem}
  \newtheorem{lemma}{Lemma}
  \newtheorem{proposition}{Proposition}
  \newtheorem{conjecture}{Conjecture}
  \theoremstyle{definition}
  \newtheorem{gap}{Gap}
  \theoremstyle{remark}
  \newtheorem{remark}{Remark}

  \title{$B_{\mathrm{base}} = N_{\mathrm{eff}}^{3/2}$ — Rigorous Motivation of the Exponent\\
    \large Why 3 and why 2, derived from $\mathcal{S}[\Theta]$ and $\mathbb{C}\otimes\mathbb{H}$}
  \author{Ing.\ David Jaroš}
  \date{2026-03-08 (v59)}

  \begin{document}
  \maketitle
\fi

\section{Approach A2: \texorpdfstring{$B_{\mathrm{base}} = N_{\mathrm{eff}}^{3/2}$}{B\_base = N\_eff\^{}3/2} — Rigorous Motivation of the Exponent}
\label{sec:hausdorff_approach}

\textbf{Overall status: [MOTIVATED CONJECTURE — with explicit gap].}
The exponent $3/2$ is \emph{algebraically derived} from the structure of
$\mathbb{C}\otimes\mathbb{H}$ and the Gaussian path integral.
Neither factor (the 3 nor the 2) is a free parameter.
What remains unproved is stated explicitly in Section~\ref{sec:gap}.
No circular use of $\alpha^{-1} = 137$ anywhere in this document.

%--------------------------------------------------------------------
\subsection{Section 1: Where does the 3 come from?}
\label{sec:three}
%--------------------------------------------------------------------

\textbf{Claim} [PROVED]: The relevant mode space for vacuum polarisation in UBT
is $\mathrm{Im}(\mathbb{H})$, which has $\dim_\mathbb{R} = 3$.

\begin{lemma}[Relevant mode space — PROVED]
\label{lem:mode_space}
  The vacuum polarisation integral in UBT runs over modes
  $\delta\Theta_V \in \mathrm{Im}(\mathbb{H})$.
  The real dimension of this space is
  \begin{equation}
    d \;:=\; \dim_{\mathbb{R}}\bigl(\mathrm{Im}(\mathbb{H})\bigr) = 3.
    \label{eq:dim_ImH}
  \end{equation}
  \begin{proof}
    The UBT field decomposes as $\Theta = \Theta_S + \Theta_V$, where
    $\Theta_S = \mathrm{Sc}(\Theta) \in \mathbb{R}$ is the scalar part and
    $\Theta_V = \mathrm{Vec}(\Theta) \in \mathrm{Im}(\mathbb{H})$ is the
    vector (pure imaginary quaternion) part.
    \textit{Why $\Theta_S$ decouples:} the vacuum polarisation tensor
    $\Pi_{\mu\nu}$ is antisymmetric by gauge invariance; the scalar part
    $\Theta_S$ is a singlet under $\mathrm{SU}(2)$ and does not carry
    gauge degrees of freedom, so it does not contribute to the running of
    the gauge coupling.
    Only $\Theta_V \in \mathrm{Im}(\mathbb{H}) = \mathrm{span}\{i,j,k\}$
    carries gauge degrees of freedom.
    The standard real basis $\{i, j, k\}$ of $\mathrm{Im}(\mathbb{H})$
    has exactly three elements; hence $d = 3$.
  \end{proof}
\end{lemma}

This is not a new assumption: the same identification
$N_{\mathrm{phases}} = d = \dim_\mathbb{R}(\mathrm{Im}\,\mathbb{H}) = 3$
already appears in the proved $N_{\mathrm{eff}} = 3 \times 2 \times 2 = 12$
counting (Theorem~N.3.1 in the $N_{\mathrm{eff}}$ derivation chain).
The lemma re-uses an already established fact.

%--------------------------------------------------------------------
\subsection{Section 2: Where does the 2 come from?}
\label{sec:two}
%--------------------------------------------------------------------

\textbf{Claim} [PROVED as a mathematical identity]: The factor of 2 in the
denominator of the exponent $d/2$ is a universal property of Gaussian integration.

\begin{proposition}[Gaussian factor — mathematical identity]
\label{prop:gaussian}
  In the one-loop path integral over fluctuations
  $\delta\Theta \in \mathrm{Im}(\mathbb{H})$,
  the Gaussian integration
  \begin{equation}
    \int \mathcal{D}(\delta\Theta)\;
      \exp\!\Bigl(-\tfrac{1}{2}\,\delta\Theta \cdot S''[\Theta_0] \cdot \delta\Theta\Bigr)
    \;=\; \bigl(\det S''[\Theta_0]\bigr)^{-1/2}
    \quad \times \;(\text{constants}),
    \label{eq:gaussian_path_integral}
  \end{equation}
  where the exponent $-1/2$ on $\det S''$ is a \emph{universal mathematical
  identity}, independent of the specific form of $\mathcal{S}[\Theta]$.
  \begin{proof}
    Step 1 (functional expansion):
    expand around the vacuum $\Theta_0$:
    \begin{equation}
      \mathcal{S}[\Theta_0 + \delta\Theta]
        \;\approx\; \mathcal{S}[\Theta_0]
        + \tfrac{1}{2}\,\delta\Theta \cdot \mathcal{S}''[\Theta_0] \cdot \delta\Theta
        + \mathcal{O}(\delta\Theta^3).
      \label{eq:action_expand}
    \end{equation}
    (The linear term vanishes because $\Theta_0$ is a saddle point.)

    Step 2 (finite-dimensional Gaussian identity):
    for any positive-definite symmetric matrix $A \in \mathbb{R}^{d\times d}$,
    \begin{equation}
      \int_{\mathbb{R}^d} \exp\!\Bigl(-\tfrac{1}{2}\,x^T A\, x\Bigr)\,d^d x
        \;=\; \frac{(2\pi)^{d/2}}{\sqrt{\det A}}
        \;=\; (2\pi)^{d/2}\,(\det A)^{-1/2}.
      \label{eq:gauss_finite}
    \end{equation}
    The exponent $-1/2$ multiplying $\log\det A$ is exact and independent of $d$.

    Step 3 (one-loop effective action):
    combining Steps 1 and 2, the partition function gives
    \begin{equation}
      \Gamma_{\mathrm{eff}}[\Theta_0]
        = \mathcal{S}[\Theta_0]
        + \tfrac{1}{2}\log\det\bigl(\mathcal{S}''[\Theta_0]\bigr)
        + \mathcal{O}(\hbar^2).
      \label{eq:Gamma_one_loop}
    \end{equation}
    The prefactor $\tfrac{1}{2}$ in front of $\log\det$ is the direct
    consequence of $(\det A)^{-1/2}$ in~\eqref{eq:gauss_finite}.
  \end{proof}
\end{proposition}

\textbf{Consequence for the renormalisation coefficient $B$:}
The one-loop contribution to $B$ from the $d$-dimensional mode space scales as
\begin{equation}
  \delta B \;\sim\; \tfrac{1}{2}\,\mathrm{Tr}\log(\mathcal{S}'')
           \;\sim\; \tfrac{d}{2}\,[\text{UV divergence per mode}],
  \label{eq:deltaB}
\end{equation}
where $d = \dim_\mathbb{R}(\mathrm{Im}\,\mathbb{H}) = 3$ from Lemma~\ref{lem:mode_space}.
The factor $1/2$ originates entirely from the Gaussian identity~\eqref{eq:gauss_finite}
and is \emph{not a free parameter}.

%--------------------------------------------------------------------
\subsection{Section 3: Combining — the exponent is \texorpdfstring{$3/2$}{3/2}}
\label{sec:exponent}
%--------------------------------------------------------------------

Combining Lemma~\ref{lem:mode_space} (the 3) and
Proposition~\ref{prop:gaussian} (the 2):

\begin{conjecture}[$B_{\mathrm{base}}$ exponent — MOTIVATED CONJECTURE with explicit gap]
\label{conj:B_exponent}
  If the one-loop renormalisation of the UBT effective potential runs over
  modes $\delta\Theta \in \mathrm{Im}(\mathbb{H})$ with
  $d = \dim_\mathbb{R}(\mathrm{Im}\,\mathbb{H}) = 3$,
  and the path-integral measure is Gaussian, then
  \begin{equation}
    B_{\mathrm{base}}
      \;=\; N_{\mathrm{eff}}^{\,d/2}
      \;=\; N_{\mathrm{eff}}^{\,\dim_\mathbb{R}(\mathrm{Im}\,\mathbb{H})/2}
      \;=\; N_{\mathrm{eff}}^{3/2}
      \;=\; 12^{3/2}
      \;=\; \sqrt{12^3}
      \;=\; \sqrt{1728}
      \;\approx\; 41.57.
    \label{eq:B_exponent_main}
  \end{equation}
  The factor $3$ comes from the algebra: $\mathrm{Im}(\mathbb{H}) = \mathrm{span}\{i,j,k\}$
  has exactly three real dimensions.
  The factor $2$ comes from mathematics: Gaussian integrals always give
  $(\det A)^{-1/2}$, i.e.\ a factor $1/2$ in the exponent.
  \emph{Neither factor is a free parameter.}
\end{conjecture}

\textbf{Dimensional consistency table.}
The formula $B_{\mathrm{base}} = N_{\mathrm{eff}}^{d/2}$ is tested across algebras:

\begin{center}
\begin{tabular}{lcccc}
  \hline
  Algebra & $d = \dim_\mathbb{R}(\mathrm{Im})$ & $N_{\mathrm{eff}}$ & $B_{\mathrm{base}} = N_{\mathrm{eff}}^{d/2}$ \\
  \hline
  $\mathbb{C}$ (QED limit) & 1 & 1  & $1^{1/2} = 1$        \\
  $\mathbb{H}$ (pure quaternion) & 3 & 3  & $3^{3/2} \approx 5.20$ \\
  $\mathbb{C}\otimes\mathbb{H}$ (UBT) & 3 & 12 & $12^{3/2} \approx 41.57$ \\
  \hline
\end{tabular}
\end{center}

The UBT row gives the phenomenologically required $B_{\mathrm{base}} \approx 41.57$.
No free parameters appear.

%--------------------------------------------------------------------
\subsection{Section 4: What remains for a full proof}
\label{sec:gap}
%--------------------------------------------------------------------

The argument above establishes \emph{why} the exponent $3/2$ is natural
and algebraically motivated.
It does \emph{not} yet constitute a full proof.
The remaining gaps are stated explicitly:

\begin{gap}[Remaining for full proof]
\label{gap:remaining}
  The derivation is complete except for two steps:
  \begin{enumerate}
    \item[(a)] \textbf{Explicit computation of $\det(\mathcal{S}'')$.}
      It must be shown that $\mathcal{S}''[\Theta_0]$ on $\mathrm{Im}(\mathbb{H})$,
      after renormalisation, gives precisely $N_{\mathrm{eff}}^{3/2}$.
      This requires an explicit computation of $\det(\mathcal{S}'')$
      (or equivalently of the functional determinant on the mode space),
      which has not yet been carried out.
      Equivalently (Remark~\ref{rem:d_eff}): the physical renormalisation
      must lift the effective dimension from
      $d_{\mathrm{eff}}(B_0) \approx 2.595$ to $d = 3.000$ ($\Delta d = 0.405$).
      Two candidate mechanisms (Seeley--DeWitt curvature correction on
      $\mathrm{Im}(\mathbb{H})$ and mode-pair coherence) were investigated in
      \texttt{b\_base\_delta\_d.tex} and found to be dead ends
      (approaches C1 and C2 both \textbf{[DEAD END]}).
      \textit{Status: \textbf{[OPEN]}}.

    \item[(b)] \textbf{Protection of the exponent at higher loops.}
      It must be shown that higher-loop corrections do not modify the exponent
      $3/2$ — i.e.\ that the $d/2$ structure is \emph{protected} by some
      symmetry or non-renormalisation theorem.
      \textit{Status: \textbf{[OPEN]}}.
  \end{enumerate}

  Until (a) and (b) are demonstrated, the overall status is:\\
  \centerline{\textbf{[MOTIVATED CONJECTURE — strong algebraic motivation,
  derivation from $\mathcal{S}[\Theta]$ not yet complete]}}
\end{gap}

\begin{remark}
  The Gaussian identity~\eqref{eq:gauss_finite} and
  Lemma~\ref{lem:mode_space} are individually \emph{proved}.
  The conjecture is in the step connecting $(\det \mathcal{S}'')^{-1/2}$
  on the mode space to the specific value $N_{\mathrm{eff}}^{3/2}$.
  The structural origin of both factors is clear; the quantitative computation
  is the open problem.
\end{remark}

\begin{remark}[Effective dimension — Approach A4 consistency check]
  \label{rem:d_eff}
  Define the \emph{effective dimension} $d_{\mathrm{eff}}(B) := 2\log B / \log N_{\mathrm{eff}}$,
  so that $B = N_{\mathrm{eff}}^{d_{\mathrm{eff}}/2}$.  Then:
  \begin{itemize}
    \item $d_{\mathrm{eff}}(B_0) = 2\log(2\pi N_{\mathrm{eff}}/3)/\log(N_{\mathrm{eff}}) \approx 2.595$
      \quad [non-integer, from the proved one-loop result];
    \item $d_{\mathrm{eff}}(B_{\mathrm{base}}) = 3.000$ exactly \quad
      [integer = $\dim_\mathbb{R}(\mathrm{Im}\,\mathbb{H})$, algebraic identity].
  \end{itemize}
  Gap~(a) is therefore equivalent to explaining why the physical renormalisation
  lifts the effective dimension from $d_{\mathrm{eff}}(B_0) \approx 2.60$ to
  the algebraic value $d = 3.00$ (a shift $\Delta d \approx 0.40$).
  Numerical details are in \texttt{tools/compute\_B\_effective\_dimension.py}
  (Approach~A4, v59).
\end{remark}

%--------------------------------------------------------------------
\subsection{Section 5: In plain language}
\label{sec:plain}
%--------------------------------------------------------------------

\textbf{In plain language:}
The UBT field $\Theta$ lives in the biquaternion algebra $\mathbb{C}\otimes\mathbb{H}$.
When we integrate out quantum fluctuations around the vacuum,
only the imaginary part $\mathrm{Im}(\mathbb{H}) = \{i, j, k\}$ contributes to
the running of the coupling — the scalar part decouples because it is not a
gauge degree of freedom.
$\mathrm{Im}(\mathbb{H})$ has exactly three real dimensions, one for each of $i$, $j$, $k$.
Gaussian integration in $d$ dimensions always produces a factor of $d/2$ in the
exponent of the result — this is a \emph{mathematical identity}
(equation~\eqref{eq:gauss_finite}), not a physical assumption.
Combining: exponent $= d/2 = 3/2$.
This is why $B_{\mathrm{base}} = N_{\mathrm{eff}}^{3/2}$ and not $N_{\mathrm{eff}}^1$
or $N_{\mathrm{eff}}^2$.
The 3 comes from the algebra of quaternions; the 2 comes from the mathematics of
Gaussian integrals.
Neither is a free parameter.

\ifdefined\INCLUDEMODE\else
\end{document}
\fi
